\section{Related Work}
\label{sec:related_work}

\paragraph{Agentic retrieval and query reformulation.}
Retrieval-augmented generation grounds model outputs in external
evidence~\cite{lewis2020retrieval}. Beyond fixed retrieve--generate pipelines,
Self-RAG learns to retrieve and critique evidence through
self-reflection~\cite{asai2024selfrag}, while Iter-RetGen alternates retrieval
and generation so that intermediate reasoning can guide subsequent
retrieval~\cite{shao2023iterretgen}. Search-R1 and R1-Searcher further use
reinforcement learning to train language models to issue search actions during
reasoning~\cite{jin2025searchr1,song2025r1searcher}. A related line focuses on
retriever-aware query reformulation. TTR generates reasoning-enhanced rewrites
for domain-specific retrievers~\cite{li2026ttr}, whereas CRAF recovers latent
diagnostic intent for clinical similar-case retrieval through reinforcement
learning~\cite{lin2026craf}. Together, these studies establish retrieval as a
sequential decision process rather than a single preprocessing step, yet their
controllers condition on the current round alone, so evidence and negative
outcomes accumulated so far are not carried forward. Our setting extends this
view from search or rewriting alone to state-carrying control that jointly
performs candidate selection and round-level decisions: one policy chooses which
evidence to retain, updates a current best evidence state, and either preserves
the current pool or redirects retrieval under a fixed budget. Retained evidence
and failed retrieval traces are represented separately so that only
answer-supporting context is exposed to the generator.

\paragraph{Multimodal and medical retrieval adaptation.}
Multimodal RAG extends external grounding beyond text. MuRAG retrieves from
image and text collections for knowledge-intensive visual question
answering~\cite{chen2022murag}; retrieval-augmented multimodal language modeling
incorporates retrieved multimodal demonstrations into generation
~\cite{yasunaga2023retrieval}; and GeMKR generates multimodal knowledge clues
before database retrieval~\cite{long2024gemkr}. More recent systems integrate
retrieval with model optimization: MMRAG-RFT applies reinforcement fine-tuning
to multimodal document filtering, ranking, and answer
generation~\cite{zhao2026mmragrft}, while URaG unifies page retrieval and
generation within a multimodal model~\cite{shi2026urag}. These works demonstrate
the complementarity of visual and textual evidence, but their discriminative
retrieval value can vary by query. We therefore treat query-dependent modality
allocation as adaptation of the fixed retrieval environment, rather than as
the central policy contribution.

Medical multimodal reasoning introduces additional domain constraints.
Biomedical vision-language pretraining and instruction tuning improve medical
knowledge and instruction following~\cite{li2023llavamed}, while EndoBench
documents substantial variation in endoscopic reasoning across clinical tasks,
scenarios, and visual prompting granularities~\cite{liu2025endobench}. CMID
further shows that medically relevant visual evidence can be obscured by
textual priors or irrelevant image regions~\cite{zhu2026cmid}. These works
motivate domain adaptation but are not themselves anatomical retrieval
methods. In endoscopic retrieval, findings from different gastrointestinal
sites may appear visually or linguistically similar while requiring different
supporting evidence. We accordingly use anatomical semantic partitioning as a
fixed medical prior in the retrieval environment and evaluate it separately
from controller learning.

\paragraph{Positioning of AgenticRL.}
AgenticRL lies at the intersection of agentic search, retriever-aware
reformulation, multimodal RAG, and medical-domain adaptation, but targets a
distinct frozen-stack control problem. During controller optimization, the
retrievers, deployed index, enabled routing modules, and generator remain
fixed. At each round, the controller jointly produces
candidate-level \textsc{Keep}/\textsc{Drop} decisions and a round-level
\textsc{Accept}/\textsc{Rewrite} action. Retained partial evidence $M_t^+$
is maintained as a current best evidence state for later rounds and the final
generator, whereas failed
retrieval traces $M_t^-$ are controller-only state used to inform subsequent
reformulation. Moreover, a rewrite is evaluated only after it has traversed
the deployed retrieval path, placing \textsc{Accept} and \textsc{Rewrite} on a
shared generator-derived utility scale. Experimentally, we compare memory-aware
control with one-shot, single-round, and memoryless iterative alternatives on
the same frozen stack, and isolate the contribution of cross-round state itself.
